\section{总结}\label{sec:summary}

在格点上计算的准光前关联的正确重正化，一直是把 LaMET 应用于部分子物理研究的主要挑战。这类关联含有与 Wilson 线相伴的线性发散，必须以高精度消除它们才能提取想要的物理信息。本文提出一种自重正化方法，从一个可在短距离处匹配到连续方案的准光前矩阵元中，同时消除线性发散、对数发散以及离散化误差。作为一个典型的范例，我们证明该方法对 $\pi$、核子和离壳夸克态中的 $O_{\gamma_t}$ 矩阵元行之有效。我们的分析表明，重正化因子在所考虑的强子态中是普适的；而且 clover 与 overlap 费米子在 $\pi$ 态给出的重正化关联彼此相近。此外，我们发现以往常用的 RI/MOM 和比值重正化方案中存在大的非微扰效应，这在最近提出的混合重正化方案中可以并且必须避免。我们的方法对 HYP 涂抹矩阵元同样有效，这无疑十分有意义，为对涂抹的高精度格点数据实施唯象重正化铺平了道路。

为便于读者查阅，我们把前几节讨论的各种矩阵元与重正化相关的参数汇集于表~\ref{tab:para_hyp0} 和表~\ref{tab:para_hyp1}。对 HYP 涂抹情形，$\Lambda_{\rm QCD}$ 的选取有较大自由度。我们在表~\ref{tab:para_hyp1} 中取 $\Lambda_{\rm QCD}=0.39$ GeV；但我们也发现 $\Lambda_{\rm QCD}=0.1$ GeV 给出更小的拟合离散化误差。


\begin{table}[htbp]
  \centering
  \begin{tabular}{c|cccc}
  \toprule
 情形 & $k$ & $\Lambda_{\rm QCD}$(GeV) & $d$ & $m_{0}$(GeV) \\
\hline
overlap 夸克 & 1.46 & 0.109(02) & -1.29 & 0.0462(16)\\
\hline
clover 夸克 & 1.46 & 0.135(02) & -1.35 & 0.1513(16)\\
\hline
overlap $\pi$ & 1.46 & 0.093(10) & -1.17 & -0.0357(46)\\
\hline
clover $\pi$ & 1.46 & 0.086(14) & -0.92 &  -0.0715(50)\\
\hline
clover 核子 & 1.46 & 0.093(06) & -0.97 &  -0.0508(40)\\
\hline
  \end{tabular}
  \caption{基于拟合函数式~(\ref{eq:logM}) 与式~(\ref{eq:gztoMS})、$\delta_{\rm sys}=0.002$ 时 MILC 与 RBC 系综无 HYP 涂抹情形的重正化参数。此处 $k$ 无量纲。}
  \label{tab:para_hyp0}
\end{table}


\begin{table}[htbp]
  \centering
  \begin{tabular}{c|cccc}
  \toprule
 情形 & $k$ & $\Lambda_{\rm QCD}$(GeV)  \\
\hline
overlap 夸克 & 0.5521(07) & 0.39 \\
\hline
clover 夸克 & 0.6328(05) & 0.39 \\
\hline
overlap $\pi$ & 0.5191(14) & 0.39 \\
\hline
clover $\pi$ & 0.5178(18) & 0.39 \\
\hline
clover 核子 & 0.5139(54) & 0.39 \\
\hline
Wilson 圈   & 0.4921(02) &  0.39 \\
\hline
VEV           & 0.39(21)  &  0.39 \\
\hline
Wilson 链接   & \multicolumn{1}{p{1cm}<{\centering}}{0.5402(04), 0.7099(09)} & 0.39  &     &  \\
\hline
  \end{tabular}
  \caption{基于拟合函数式~(\ref{eq:logM_ss}) 的 MILC 与 RBC 系综 HYP 涂抹情形的重正化参数。$k$ 为 $e(z)$ 的斜率（Wilson 圈取斜率之半），此处无量纲。}
  \label{tab:para_hyp1}
\end{table}


出于比较的目的，我们还展示了基于 CLS 合作组 $2+1$ 味 clover 夸克与 Luescher-Weisz（等价于 Symanzik）规范系综的结果~\cite{Bruno:2014jqa}。价费米子作用量取为 clover 作用量，与海费米子作用量相同。CLS 系综的重正化参数见表~\ref{tab:CLS}。CLS 系综 $\pi$ 矩阵元的参数 $k$ 与 $\Lambda_{\rm QCD}$ 同 MILC 与 RBC 系综的结果相近，提示无论是否采用混合作用量，结果都是一致的。

\begin{table}[htbp]
  \centering
  \begin{tabular}{c|cccc}
  \toprule
 情形 & $k$ & $\Lambda_{\rm QCD}$(GeV) & $d$ & $m_{0}$(GeV) \\
\hline
clover 夸克 hyp0 & 1.46 & 0.170(06)  &  -0.68 &  0.1936(26) \\
\hline
clover $\pi$ hyp0 & 1.46 & 0.113(09)  &  -2.24  &  -0.0941(49) \\
\hline
clover 夸克 hyp1 & 0.573(05)  & 0.39 & 0.53  &  0.2298(26) \\
\hline
clover $\pi$ hyp1 & 0.486(10) & 0.39 & -0.21 &  0.0312(29) \\
\hline
  \end{tabular}
  \caption{$\delta_{\rm sys}=0.002$ 时基于拟合函数式~(\ref{eq:logM}) 与式~(\ref{eq:gztoMS}) 的 CLS 系综重正化参数。“hyp0”为未涂抹情形，“hyp1”为 HYP 涂抹情形。此处 $k$ 无量纲。}
  \label{tab:CLS}
\end{table}

我们想强调的是，本方法的应用并不限于本文所讨论的这些矩阵元。原则上，它可以应用于式~(\ref{eq:operator}) 形式的任何算符矩阵元，
也可以推广到 LaMET 应用中含有 Wilson 线的任意矩阵元~\cite{Ji:2020ect}。当然，在这些情形中还需要仔细考虑矩阵元的物理解释以及与格点计算相关的细节。
